% ! TeX root = distributed-systems-final-report.tex
\section{Validation}\label{validation}

\subsection{Automatic Testing}

In the development of this project, an effort was made initially to do
\textbf{unit-testing} in a \textbf{test-driven} approach, as stated in the
requirement \ref{req:test}, as a personal new experience. However, the promise
wasn't really maintained.

Only in the very first stage, during the development of the first version of the
simulated world, without any controller or vehicle logic but only with physical
movements and mathematical operations, every time that a functionality was
needed, a test was developed first and only then that would have been
implemented with a separate function (so that it could be tested alone).
Unfortunately, this didn't last long, as the required features, and the related
functions, soon became too complex. In the first place, this means that a
function becomes too hard to test, with too many possible normal or edge cases,
and unexpected ones. The main problem, however, is that it was hard to even know
in advance how the controller's or vehicle's logic should have come out at the
end, so developing a test for a function that would have been changed or even
removed soon after would have taken too much time away from the distributed
systems related topics.

That said, just a few unit tests were developed right at the end, to test only a
few isolated functionalities.
\\
\\
Proper \textbf{end-to-end} test and \textbf{deployment automation} were never
even taken in consideration at all, given the lack of time, the total absence of
personal experience with specific tools and the probably complex general
architecture.

For communication and interaction testing, manual tests was deemed enough
(\ref{validation-acceptance}).
\\
\\
\textit{pytest} (9.1.1) has been used for unit testing. All tests (in the
$src/tests/$ folder) can be run, from the project's root directory, with the
command:
\begin{verbatim}
  pytest
\end{verbatim}

One should first install the packages required by the tested services (i.e. all
except the web viewer), besides pytest itself. $src/tests/requirements.txt$
contains them, and can be used with:
\begin{verbatim}
  pip install -r src/tests/requirements.txt
\end{verbatim}

The developed unit tests for physics ($test\_physics.py$) and mathematics
($test\_math\_utils.py$ and $test\_roundabout.py$) regard kinematics and
calculations: they test, for a few cases, that distances are calculated
correctly, for example on a line or on a circle, and that speeds and
accelerations (positive or negative) are updated and used correctly, i.e.
vehicles ride for the expected distance or the expected amount of time.

Unit tests for the the controller ($test\_controller.py$) test the command
evaluation in a few cases by creating mock scenarios, verifying that a
tail-gating vehicle decelerates to increase the margin from the one on the
front, or that a vehicle with a lot of margin on the front keeps accelerating.

Finally, unit tests for the vehicle ($test\_vehicle.py$) perform analogous tests
to those of the controller, but using its failsafe mode logic, and also check
the collision detection.

\subsection{Acceptance test}\label{validation-acceptance}

Given the lack of automatic testing, manual testing has clearly been fundamental
in the project, to continuously debug it and eventually validate it.

Being the subject a simulated real-world scenario, verifying that everything
goes as expected by looking at the visualization is straight-forward, although
debugging requires much more effort, by deeply analyzing the code and
correlating it with what happens in the simulation, or by checking logs and
vehicles and controller statistics. The behaviour one should expect to see is
what described in the requirements (\ref{req:requirements}).

For communication and interaction between the services, the simple fact that the
project worked was enough: bugs and problems only came after, when the it was
the logic to be put under examination.
\\
\\
The developed system commands to force the state of vehicles to failsafe or
disconnected mode allowed to test these other failure scenarios.

The ability to pause the simulation was also fundamental, to either stop to
check logs and statistics, or to watch from up close the interaction between two
cars. Moreover, by pausing, one could change the state of vehicles with
precision, for example to set to disconnected one which is already in a
dangerous situation, in order to test edge cases.
\\
\\
The simulation has thus been run with different numbers of vehicles (between 10
and 50), to see how the system and the controller's work would change. However,
the road configuration has never been changed, although it's possible to easily
change the number of roads or the size of the roundabout.

Once run the simulation, tests were made with varying amounts of vehicles into
failsafe of disconnected mode.
\\
\\
Some possible evaluation metrics at disposal, available on the simulation page
(refer to the user guide, \ref{user-guide}), are the number of happened
collisions and the total number of spawned vehicles in a certain amount of time,
to measure the total throughput of the roundabout. Here follow some example
runs; refer to the self-evaluation (\ref{self-evaluation}) for a brief analysis.

Each test has been run for 10 minutes:
\begin{itemize}
  \item 40 vehicles (all in normal state): 309 spawned, 9 collisions
  \item 40 vehicles (all in normal state): 325 spawned, 7 collisions
  \item 40 vehicles (30 in failsafe): 335 spawned, 49 collisions

  \item 40 vehicles (all disconnected): 288 spawned, collisions unavailable (the
  system went in roundabout deadlock after 8 minutes)
  \item 40 vehicles (all disconnected): 352 spawned, 72 collisions
  \item 40 vehicles (10 disconnected): 304 spawned, 37 collisions

  \item 20 vehicles (all in normal state): 360 spawned, 23 collisions
  \item 10 vehicles (all in normal state): 221 spawned, 3 collisions
\end{itemize}
